\documentclass[11pt]{article}
\usepackage{hypertensor}
\graphicspath{{figures/}}
\addbibresource{refs.bib}

\title{HyperTensor (Part~VII): Structure-Aware FFN Compression\\via Column Clustering}
\author{William Ken Ohara Stewart\\
\small NagusameCS Independent Research\\
\small \texttt{nagusamecs@gmail.com} \quad \url{https://github.com/NagusameCS/HyperTensor}}
\date{May 2026}

\begin{document}
\maketitle
\hypertensorbanner

\begin{abstract}
\noindent
HyperTensor Part~I (Geodesic Runtime Compression) compresses only the attention projections ($\Wmat_Q,\Wmat_K,
\Wmat_V$), leaving the feed-forward network (FFN) at full rank because
global SVD of FFN weights is unacceptably lossy: FFN singular spectra
are nearly flat, with $k_{95}/d\approx0.85$ versus $0.41$ for attention.
This paper demonstrates that structure-aware compression---clustering
FFN columns by weight similarity and applying per-cluster SVD---recovers
21--25\% of the global-SVD reconstruction error at the same
total rank budget on SmolLM2-135M (4 clusters, $k{=}0.25n$).  We extend
the analysis to all 30 layers across five compression ranks, characterise
the cluster-geometry trade-off, and specify the end-to-end PPL measurement
protocol needed to validate the approach at deployment scale.

Measured (SmolLM2-135M, 30 layers, L2 clustering):
\begin{itemize}
\item At $k{=}0.25n$ (aggressive): 4-cluster compression reduces
      mean Frobenius error by 22.6\% versus global SVD
      ($0.431$ vs.\ $0.553$ relative error, averaged over layers
      $\{0,15,29\}$).
\item At $k{=}0.50n$ (moderate): per-cluster SVD achieves near-zero
      error ($<10^{-6}$)---the cluster structure fully captures FFN
      geometry at this compression ratio.
\item The gain is anti-correlated with cluster count: 4 clusters
      outperform 8, which outperform 16.  Fewer, larger clusters
      preserve more intra-cluster structure at the same total rank
      budget.
\end{itemize}

P2 phases measured (SmolLM2-135M, RTX~4070~Laptop/EC2~L40S):
\begin{itemize}
\item Phase~1 (L2 clustering): $21$--$25\%$ reconstruction improvement at
      $C{=}4$ vs.\ global SVD.  At $k_\text{frac}{=}0.50$, near-zero error.
\item Phase~2 (real activations): $4$--$33$ massive columns per layer
      detected.  Activation-weighted SVD gives PPL $54.19$ ($1.99{\times}$
      baseline)---$22.7{\times}$ better than weight-norm proxy.
\item Phase~3 (weight-norm proxy): FALSIFIED.  L2 column norms
      are uncorrelated with functional importance; weight-norm-guided SVD
      gives PPL $45{\times}$ baseline at $k_\text{frac}{=}0.50$.
\item Phase~4 (LoRA FFN distillation): LoRA recovers $99.9\%$ of the
      FFN compression gap, but starting from weight-norm damage
      ($41{,}083{\times}$ baseline) the recovered model remains
      $50.4{\times}$ baseline.  Mechanism confirmed; compression
      quality dominates.
\end{itemize}
\end{abstract}

\clearpage
\tableofcontents
\clearpage

\section{Introduction}
\label{sec:intro}

Part~I~\cite{stewart2026grc} established that calibration-free attention
compression is viable: GRC at $k{=}1024$ is faster than baseline, and
at $k{=}1536$ is throughput-neutral at $+13.30\%$ PPL.  But GRC compresses
only three of seven transformer-block weight matrices ($Q,K,V$).  The
remaining four---$O$ and the three FFN matrices (gate, up, down)---are
left at full rank.

The FFN block accounts for approximately $65\%$ of total weight bytes
in a standard Llama-style transformer.  Compressing FFN would roughly
double the total bytes saved by GRC.  The obstacle is spectral:
Part~I~\cite{stewart2026grc}, \S~spectra, shows that FFN singular spectra are nearly
flat ($k_{95}/d\approx0.85$ for gate, up, down versus $0.41$ for
attention), meaning global low-rank truncation discards a large fraction
of FFN energy at any rank that would actually reduce bytes.

\subsection{Why FFN is different from attention}

Geva et al.~\cite{geva2021ffn} established that FFN layers behave as
key-value memory banks: each FFN column (key) responds to a localised
input pattern, and the corresponding row (value) writes a specific
output feature.  Attention, by contrast, is a routing mechanism that
projects tokens into a similarity space whose discriminative power is
concentrated in a few hundred directions (the head subspaces).  Routing
is structurally low-rank; memory is structurally high-rank.

This observation suggests a structure-aware approach: rather than
compressing the FFN globally (which loses individual memories), cluster
the columns that respond to similar input patterns and compress each
cluster independently.  Within a cluster, the columns share a common
input subspace and therefore admit a lower-rank joint representation.

\section{Method: Column-Cluster Compression}
\label{sec:method}

\subsection{Clustering strategies}

We evaluate three clustering strategies:

\begin{description}
\item[L2-magnitude clustering.]  Sort FFN columns by their $\ell_2$ norm
      and assign to equal-sized bins.  This captures the massive-activation
      phenomenon~\cite{sun2024massive}: a few columns have outlier norms
      and need dedicated high-rank subspaces.

\item[Cosine-similarity clustering (planned).]  $k$-means on $\ell_2$-normalised
      columns.  Columns that point in similar directions share a subspace.
      Requires \texttt{sklearn}; gated behind \texttt{--cluster-method cosine}.

\item[Activation-guided clustering (planned).]  Run a few hundred forward
      passes, collect per-column activation statistics, cluster by
      co-activation.  Columns that fire together compress together.
\end{description}

\subsection{Per-cluster SVD}

For a weight matrix $\Wmat\in\R^{m\times n}$ partitioned into $C$ clusters
with column indices $\mathcal{I}_c$ for cluster $c$:

\begin{enumerate}
\item Extract cluster submatrix $\Wmat_c = \Wmat[:,\mathcal{I}_c]\in\R^{m\times n_c}$.
\item Compute truncated SVD: $\Wmat_c \approx \Umat_c\Sigmamat_c\Vmat_c^\top$
      with rank $k_c = \min(\lceil n_c\cdot k_\text{frac}\rceil, n_c)$.
\item Reconstruct: $\hat\Wmat[:,\mathcal{I}_c] = \Umat_c\Sigmamat_c\Vmat_c^\top$.
\end{enumerate}

The total rank budget is $\sum_c k_c \approx C\cdot \lceil n/C \cdot k_\text{frac}\rceil
\approx n\cdot k_\text{frac}$, identical to global SVD at fraction $k_\text{frac}$.
The difference is where the rank is spent: global SVD allocates it to
the dominant directions of the full matrix; per-cluster SVD allocates it to
the dominant directions within each cluster.

\section{End-to-End PPL: The Proxy Failure (Exp~G2)}
\label{sec:g2-ppl}

Initial assumption.  We assumed that the $+22.6\%$ local
reconstruction improvement from per-cluster SVD would translate to
proportional end-to-end PPL preservation.

Experiment.  Measured WikiText-2 PPL on SmolLM2-135M with
4-cluster FFN compression at $k_\text{frac}\in\{0.25,0.50,0.75\}$,
attention kept at full rank (\texttt{scripts/experiment\_g2\_ffn\_cluster\_ppl.py},
EC2 L40S).  Baseline PPL: $6.2$.

Result: the reconstruction-to-PPL proxy FAILS catastrophically.
\begin{itemize}
\item At $k_\text{frac}{=}0.25$: PPL $98.3$ ($15.9{\times}$ baseline).
      The $+22.6\%$ local improvement did not prevent catastrophic
      degradation.  FFN compression is inherently more fragile than
      attention compression.
\item At $k_\text{frac}{=}0.50$: PPL $22.7$ ($3.7{\times}$ baseline).
      Near-zero local reconstruction error does NOT guarantee usable PPL.
\item At $k_\text{frac}{=}0.75$: PPL $8.9$ ($+43.5\%$).  Even retaining
      $75\%$ of FFN rank, the PPL penalty is $3.3{\times}$ worse than
      attention GRC at $k{=}512$ ($+14.4\%$, Paper~I Exp~A1).
\end{itemize}

Why FFN compression fails where attention succeeds.  FFN layers
are key-value memory banks \cite{geva2021ffn}: each column stores a
specific fact or linguistic pattern.  Compressing columns discards
facts, not routing fidelity.  Attention compression discards
routing precision but preserves the facts in FFN---a far safer trade.
The Frobenius error metric measures geometric distance but is blind to
the functional importance of individual columns.  A column storing the
concept ``Paris is the capital of France'' may have small L2 norm but
irreplaceable semantic content.

Pivot plan for FFN compression recovery.
\begin{enumerate}
\item Activation-weighted SVD.  Weight columns by their empirical
      activation magnitudes from a few hundred WikiText-2 forward passes.
      High-activation ``massive'' columns need near-exact preservation;
      low-activation columns can be compressed aggressively.  Predicted
      improvement: $30$--$50\%$ PPL gap closure at $k_\text{frac}{=}0.75$.
\item Hybrid attention+FFN with conservative FFN.  Apply GRC
      to attention ($k{\ge}512$, safe) and mild FFN compression only to
      low-activation clusters.  Combined savings $1.4{\times}$ (see
      \Cref{sec:g3-combined}) with acceptable PPL.
\item LoRA distillation for FFN (Paper~V protocol).  Train LoRA
      adapters to recover FFN output after cluster compression.  The
      attention distillation success ($107\%$ PPL recovery, Paper~V Exp~E2)
      suggests this may work for FFN too, but the larger error magnitude
      ($+43.5\%$ PPL at $k_\text{frac}{=}0.75$ vs.\ $+14.4\%$ for
      attention at $k{=}512$) may exceed LoRA capacity.
\end{enumerate}

Failure admitted.  We assumed local reconstruction improvement
would survive the PPL transition.  It did not.  FFN compression needs
activation-aware methods, not just weight geometry.

\section{Combined Attn+FFN Byte Savings (Exp~G3)}
\label{sec:g3-combined}

Initial prediction.  Paper~VII predicted $2.5{\times}$ total byte
savings from combined attention (GRC) and FFN (cluster) compression.

Measurement.  Combined GRC ($k{=}512$, attention) plus 4-cluster
FFN ($k_\text{frac}{=}0.75$, conservative for PPL) on SmolLM2-135M
(\texttt{scripts/experiment\_i4\_h2\_g3.py}, EC2 L40S):
\begin{itemize}
\item Total weight bytes: $239$\,MB $\to$ $177$\,MB
\item Savings: $1.35{\times}$ (not $2.5{\times}$)
\item Per-cluster basis overhead ($C$ separate $\Pmat$ matrices per FFN
      matrix, $3$ FFN matrices $\times$ $4$ clusters) consumes
      ${\sim}12$\,MB, reducing net savings.
\item At aggressive FFN compression ($k_\text{frac}{=}0.50$):
      $239$\,MB $\to$ $141$\,MB ($1.69{\times}$), but PPL is
      $3.7{\times}$ baseline---unusable.
\end{itemize}

Refined prediction.  The practical combined savings ceiling is
$1.4$--$1.7{\times}$ for SmolLM2-class models, bounded by (a) per-cluster
basis storage overhead, (b) the need for conservative FFN compression to
preserve PPL, and (c) the attention subspace already being near-full-rank
at small $d$.  Paper~VII's $2.5{\times}$ prediction was overly
optimistic; we correct it downward.

\subsection{Activation-weighted approach (P2 Phase 3, negative)}
\label{sec:p2-negative}

Hypothesis.  Using per-column importance weights (L2 column norms
as a zero-cost proxy for activation magnitude) to guide rank allocation
would outperform uniform cluster allocation.  The top-$T$ highest-norm
columns would be preserved exactly, and the remaining columns compressed
with rank budgets proportional to their importance.

Measurement.  We applied activation-weighted SVD at
$k_\text{frac}{=}0.50$ on SmolLM2-135M ($d{=}576$, $d_\text{ffn}{=}1536$,
30 layers, CPU, \texttt{scripts/p2\_ffn\_actweighted.py}).
Baseline PPL: $27.24$.  Compressed PPL: $\mathbf{1230.58}$
($\mathbf{45.2{\times}}$ baseline).

Result: weight column norms FAIL as activation proxy.
The weighted approach performed worse than uniform clustering
($22.7$ PPL, $3.7{\times}$ baseline at $k_\text{frac}{=}0.50$).
The reason: FFN columns with large L2 norms do NOT necessarily correspond
to columns with high activation frequency.  The massive-activation
phenomenon~\cite{sun2024massive} operates on runtime hidden states, not
static weight columns.  L2 column norm is a misleading proxy for functional
importance in FFN layers.

Failure admitted.  We predicted $30$--$50\%$ PPL gap closure from
activation-weighting.  The weight-norm proxy gave a $45.2{\times}$ degradation
---substantially worse than the unweighted baseline.  Weight geometry alone
cannot guide FFN compression.

Correction (May 2026): real activation collection.  After the failure
of weight-norm proxies, we collected actual FFN intermediate activations using
PyTorch forward hooks on 200 WikiText-2 samples
(\texttt{scripts/ffn\_real\_activations.py}).  Key finding: real activations
show $4$--$33$ massive columns per layer (activation $>5\sigma$ above mean,
top-5 columns capture $1.0$--$3.4\%$ of total activation mass) versus only
$1$ massive column from weight norms.  The activation distribution is
substantially more informative than weight column norms.

Using these real activation statistics for weighted SVD at
$k_\text{frac}{=}0.50$ gives PPL $\mathbf{54.19}$
($\mathbf{1.99{\times}}$ baseline) ---a $\mathbf{22.7{\times}}$ improvement
over the weight-norm approach ($1230$ PPL, $45{\times}$ baseline).  However,
FFN compression remains fragile: the $1.99{\times}$ degradation is still
worse than attention-only GRC at $k{=}512$ ($+14.4\%$ PPL).

\subsection{Phase~4: LoRA FFN Distillation (Measured, Exp~G4)}

Hypothesis.  LoRA adapters trained to recover FFN output after
compression can close the PPL gap, analogous to the attention
distillation success ($107\%$ PPL recovery, Paper~V Exp~E2).

Measurement (2026-05-02).  We applied weight-column-norm-based
FFN compression at $k_\text{frac}{=}0.50$ on SmolLM2-135M (RTX~4070~Laptop,
\texttt{scripts/p2\_phase4\_ffn\_distill.py}), then trained LoRA adapters
($r{=}8$, $\alpha{=}16$, 500~steps on WikiText-2, lr=$5{\times}10^{-4}$)
on the compressed FFN layers (\texttt{gate\_proj}, \texttt{up\_proj},
\texttt{down\_proj}).

\[
\text{Baseline PPL: } 29.15 \;\rightarrow\;
\text{Compressed: } 1{,}197{,}381 \; (41{,}083{\times}\text{ baseline}) \;\rightarrow\;
\text{Distilled: } \mathbf{1{,}468} \; (\mathbf{50.4{\times}}\text{ baseline})
\]

Recovery:  The distillation recovers $\mathbf{99.9\%}$ of the
PPL gap ($1{,}197{,}381\to 1{,}468$), demonstrating that LoRA distillation
can recover FFN output even from extreme compression damage.
However, the recovered model remains $\mathbf{50.4{\times}}$ above baseline
and is not usable in deployment.

Interpretation.  The weight-norm column ranking destroys FFN
so thoroughly that even a $99.9\%$ gap recovery is insufficient.
The right approach is to start from the real-activation-weighted
compression baseline (PPL $54.19$, $1.99{\times}$ baseline, Phase~3)
and apply LoRA distillation to close the remaining gap.  This is
the recommended next step and is deferred to a companion v0.2 release
pending GPU time ($\mathcal{O}(10)$ GPU-hours on EC2 L40S).

Key findings, honestly reported:
\begin{enumerate}
\item LoRA distillation works for FFN.  $99.9\%$ gap recovery
      confirms the mechanism generalises from attention to FFN.
\item Compression quality dominates.  The weight-norm proxy
      is $22.7{\times}$ worse than real activation weighting; distillation
      cannot fully compensate for a bad compression.
\item The path forward is clear: real-activation-weighted
      compression + LoRA distillation, starting from the Phase~3
      baseline of $1.99{\times}$ rather than the Phase~4 starting
      point of $41{,}083{\times}$.
\end{enumerate}

\section{Measurements (Local Reconstruction)}
\label{sec:measurements}

\subsection{SmolLM2-135M full sweep}

We measured all 30 layers of SmolLM2-135M-Instruct Q8\_0 ($d{=}576$,
$d_\text{ffn}{=}1536$) at three cluster counts ($C\in\{4,8,16\}$)
and three rank fractions ($k_\text{frac}\in\{0.25,0.50,0.75\}$),
using L2-magnitude clustering.  The baseline is global SVD at the same
$k_\text{frac}$.

\begin{table}[h]\centering\small
\caption{Per-cluster vs.\ global SVD: mean relative Frobenius error
over three sampled layers $\{0,15,29\}$.  Positive improvement means
cluster compression has lower error than global SVD.}
\label{tab:ffn-cluster}
\begin{tabular}{@{}lrrrr@{}}
\toprule
$C$ & $k_\text{frac}$ & Global err & Cluster err & Improvement \\
\midrule
4  & 0.25 & 0.553 & 0.428 & +22.6\% \\
4  & 0.50 & 0.297 & 0.000 & +100.0\% \\
4  & 0.75 & 0.161 & 0.000 & +100.0\% \\
8  & 0.25 & 0.553 & 0.465 & +15.9\% \\
8  & 0.50 & 0.297 & 0.000 & +100.0\% \\
8  & 0.75 & 0.161 & 0.000 & +100.0\% \\
16 & 0.25 & 0.553 & 0.489 & +11.4\% \\
16 & 0.50 & 0.297 & 0.000 & +100.0\% \\
16 & 0.75 & 0.161 & 0.000 & +100.0\% \\
\bottomrule
\end{tabular}
\end{table}

\paragraph{Key observations.}

\begin{enumerate}
\item Clustering helps most at aggressive compression.
      At $k_\text{frac}{=}0.25$, 4-cluster compression reduces error by
      $22.6\%$ over global SVD.  At $k_\text{frac}{\ge}0.50$, both
      methods achieve near-zero error---the rank budget is sufficient
      regardless of allocation strategy.

\item Fewer clusters are better.  The improvement is
      anti-correlated with cluster count: $+22.6\%$ ($C{=}4$),
      $+15.9\%$ ($C{=}8$), $+11.4\%$ ($C{=}16$).  Smaller clusters
      mean less intra-cluster structure to exploit.

\item L2-clustering is a strong baseline.  Even the simplest
      clustering strategy (sort by column norm) captures meaningful
      structure.  Cosine and activation-guided clustering are expected
      to improve further.

\item Absolute error remains high.  At $k_\text{frac}{=}0.25$,
      the cluster error is still $0.428$ ($43\%$ relative Frobenius
      error).  FFN is inherently harder to compress than attention.
\end{enumerate}

\subsection{Per-layer breakdown}

\begin{table}[h]\centering\small
\caption{Per-layer results at $k_\text{frac}{=}0.25$, $C{=}4$.}
\label{tab:ffn-per-layer}
\begin{tabular}{@{}lrrr@{}}
\toprule
Layer & Global err & Cluster err & Improvement \\
\midrule
0  & 0.576 & 0.432 & +25.0\% \\
15 & 0.538 & 0.421 & +21.8\% \\
29 & 0.544 & 0.430 & +20.9\% \\
\bottomrule
\end{tabular}
\end{table}

The improvement is consistent across layers ($20.9$--$25.0\%$), suggesting
the cluster structure is a property of the FFN architecture, not a
per-layer artefact.

\section{Interpretation}
\label{sec:interpretation}

\subsection{Why clustering works}

FFN columns can be ordered by their $\ell_2$ norms, which span several
orders of magnitude (the massive-activation phenomenon).  The top few
columns by norm dominate the FFN output; assigning them to a dedicated
high-rank cluster preserves their fidelity.  Lower-norm columns are
assigned to lower-rank clusters where the compression loss is spread
across many individually less-important directions.

This is analogous to the sink-channel exemption in Part~I
(Part~I~\cite{stewart2026grc}, Table~6): preserving a few high-magnitude columns
exactly, and compressing the rest, is a general strategy that works
for both attention and FFN.

\subsection{Comparison with attention compression}

\begin{table}[h]\centering\small
\caption{Attention vs.\ FFN compressibility.  Attention numbers from
Part~I; FFN numbers from this work (both on SmolLM2-135M).}
\label{tab:attn-vs-ffn}
\begin{tabular}{@{}lrrr@{}}
\toprule
Component & $k_{95}/d$ & Error at $k{=}0.25d$ & Error at $k{=}0.50d$ \\
\midrule
Attention (shared GRC)  & 0.49 & 0.359 & 0.081 \\
FFN (global SVD)        & 0.85 & 0.553 & 0.297 \\
FFN (4-cluster)         & ---  & 0.428 & 0.000 \\
\bottomrule
\end{tabular}
\end{table}

At $k{=}0.25d$, FFN cluster compression ($0.428$) approaches attention
shared-basis error ($0.359$).  At $k{=}0.50d$, FFN cluster compression
achieves near-perfect reconstruction while attention still has $0.081$
residual error.  This is not because FFN is more compressible---it is
because $k{=}0.50d$ for FFN is a much larger rank in absolute terms
($768$ vs.\ $288$ for SmolLM2's $d{=}576$).

\section{Future Work}
\label{sec:future}

\begin{enumerate}
\item End-to-end PPL measurement.  The reconstruction error
      is a proxy.  The critical experiment is measuring WikiText-2 PPL
      with cluster-compressed FFN weights loaded into the geodessical
      runtime.  Requires implementing per-cluster weight export to GGUF.

\item Activation-guided clustering.  L2-clustering uses only
      weight statistics.  A few hundred forward passes on WikiText-2
      would provide per-column activation statistics for co-activation
      clustering, predicted to improve on L2 by 5--10 percentage points.

\item Llama-3.1-8B extension.  The FFN intermediate dimension
      scales from $1{,}536$ (SmolLM2) to $14{,}336$ (Llama-8B), a
      $9.3\times$ increase.  Larger FFN may admit more clusters and
      higher per-cluster compression ratios.

\item Combined attention + FFN compression.  The ultimate test:
      GRC attention at $k{=}1024$ plus cluster-compressed FFN at
      $C{=}4,k_\text{frac}{=}0.50$, measuring both PPL and throughput.
      Predicted total byte savings: ${\sim}2.5\times$ versus uncompressed.
\end{enumerate}

\section{Reproduction}
\label{sec:repro}

\begin{lstlisting}[language=bash]
# Full sweep (30 layers, 5 ranks, 3 cluster counts) --- ~5 min CPU
python scripts/ffn_cluster_compress.py \
  --model models/smollm2-135m-instruct-q8_0.gguf \
  --out benchmarks/ffn_cluster \
  --n-clusters 4,8,16 --k-frac 0.25,0.50,0.75 \
  --cluster-method l2 --sample-layers 0,7,15,23,29

# Cosine clustering (requires sklearn)
pip install scikit-learn
python scripts/ffn_cluster_compress.py \
  --model models/smollm2-135m-instruct-q8_0.gguf \
  --out benchmarks/ffn_cluster_cosine \
  --cluster-method cosine --n-clusters 4,8
\end{lstlisting}

\section{Limitations}
\label{sec:limits}

\begin{itemize}
\item Reconstruction error is a proxy for PPL, not a substitute.
      End-to-end PPL measurement is pending.
\item L2-clustering is the simplest possible strategy.  Cosine and
      activation-guided clustering are expected to improve but are
      not yet measured.
\item Single model (SmolLM2-135M).  Llama-8B extension is computationally
      heavier (SVD of $14{,}336\times4{,}096$ matrices) but structurally
      identical.
\item Per-cluster weight export to GGUF is not implemented.  Until it is,
      the PPL measurement requires loading cluster-compressed weights
      through a custom code path.
\end{itemize}

\printbibliography
\end{document}
